\section{引言}

建筑能耗受气候条件、围护结构、室内得热和空调系统形式的共同影响，因此整栋建筑能耗模拟已经成为方案阶段进行技术比较的重要手段。EnergyPlus能够在统一框架下描述建筑负荷、系统运行控制以及冷热源侧耦合过程，因此被广泛用于建筑环境与设备工程领域的分析研究\cite{crawley2001energyplus,energyplus232docs}。与此同时，参考建筑原型方法为多气候地区下的横向比较提供了统一的几何和运行边界，使不同地区的方案差异能够在一致假定下得到识别\cite{deru2011reference,field2010reference}。

对我国公共建筑而言，气候分区差异尤其关键。严寒地区、寒冷地区、夏热冬冷地区与夏热冬暖地区在全年冷热需求构成上存在显著区别，空调技术方案的适宜性也不应预设结论，而应由负荷水平、能耗指标和设备容量共同决定。GB~50189-2015为公共建筑节能设计提供了统一规范背景，也为不同气候分区下的方案讨论提供了评价依据\cite{gb50189}。然而，传统课程设计往往以单一城市或单一系统为对象，难以在同一建模边界下完成不同气候区域和不同空调方案的平行比较。

基于此，本文构建了一个共享办公建筑原型的比较框架，并将其扩展到沈阳、天津、成都和深圳四个代表性城市。研究中，VRF+DOAS与FCU+DOAS被视为完全并列的候选方案，不预设主次。为避免系统效率掩盖气候差异，本文首先建立理想负荷基线，用以识别城市气象条件和设计日条件对建筑本体冷热需求的影响。在当前版本中，四个城市分支保持相同的建筑几何、围护结构、热分区和运行条件，从而更清晰地凸显气候因素和系统拓扑差异。所采用的逐时气象数据来自CSWD天气数据包，可保证四个城市在气象来源上的一致性\cite{onebuilding}。

本文的工作可概括为三点。第一，建立了一套可复现的办公建筑空调比较流程，用于衔接气象准备、城市分支构建、系统分支构建、结果整理和图表输出。第二，形成了覆盖四个城市、两类实际系统及理想负荷基线的完整比较矩阵。第三，明确指出当前FCU+DOAS指标的解释边界，即现阶段统计的是空调电耗而非完整的一次能源或全燃料消耗，因此相关结论应在该口径下审慎理解。
